\section{ABSTRACT}

This work assessed the molecular presence and genomic characterization of SARS-CoV-2, enterovirus and respiratory syncytial virus (RSV) in wastewater from the Quitumbe Wastewater Treatment Plant (WWTP) in Quito, in order to address the limited availability of multipathogen information in the country's wastewater. Samples were collected with a torpedo-type passive sampling device over 56 weeks and concentrated with polyethylene glycol before RNA extraction. SARS-CoV-2 was detected by RT-qPCR and its genome was amplified with the ARTIC V3 and VarSkip primer schemes; RSV amplification was attempted with an overlapping-amplicon scheme specific to RSV-A and RSV-B, run in duplicate on diluted and undiluted RNA, and enteroviruses with a nested PCR of the VP1 region. Amplicons were sequenced on Oxford Nanopore and analyzed with EPI2ME, Nextclade and Freyja for SARS-CoV-2, and by VP1 identity and phylogeny for enteroviruses. RT-qPCR detected the SARS-CoV-2 N1 gene in 31 of 42 weeks (73.8\%; 95\% CI: 58.9--84.7\%), but only six of the 21 sequenced samples, from weeks 3 to 9, yielded partial consensus sequences, assigned to JN.1-descendant lineages of clades 24A and 24H and dominated by QA.4 according to Freyja; the remaining 15 yielded no viral reads. In two of the 22 samples analyzed for enterovirus, the VP1 sequence identified Coxsackievirus A1, and no poliovirus was detected; the four samples sequenced for RSV yielded no viral reads. Retrospective samples preserved in DNA/RNA Shield 2X and in PBS 1X yielded no sequence. The results show that a single sampling and processing workflow recovers SARS-CoV-2 and a non-polio enterovirus at an urban sentinel site, and that low viral load, non-specific amplification and sample preservation limited genomic characterization. Early sample processing, standard-curve quantification and negative controls during sequencing are recommended to strengthen environmental surveillance in Quito.

\vspace{0.5\baselineskip}

\noindent\textbf{Key words:} wastewater-based surveillance, SARS-CoV-2, enterovirus, poliovirus, respiratory syncytial virus, Oxford Nanopore sequencing, genomic epidemiology, Quitumbe WWTP.
